De strijd was ongelijk.
Emmy, energierijk,
won simpel en glansrijk.

Emmy was veel sterker
als harde landwerker.
Dat bleek in de kerker,

waar ze een tijd vast zat,
omdat ze gevloekt had
en een brood had gejat,

op grond van de aanklacht
die de landheer aanbracht
als wettelijke macht.

Tussen al die slechten
liet ze zich niet knechten
en leerde te vechten.

Destijds stonden wachters
bekend als verkrachters
van de overnachters

in een gevangenis.
Tijdens haar hechtenis
ervoer ze geen crisis

wat betreft haar zeden.
De bewakers deden
zonder klare reden

juist heel welgemanierd.
Er werd wel feestgevierd,
maar zij werd niet ontsierd.

Haar aantijging bleek vals.
De landheer als jakhals
gebruikte het als wals

over haar moeder heen.
Die murw bij hem verscheen,
maanden nadat er geen

vooruitgang leek te zijn.
Het juridisch domein
kende toen een tijdlijn

die de adel diende.
Hoewel de vijftiende
dus anders verdiende

bleef Emmy’s vrijspraak uit.
Haar ma nam het besluit
te worden uitgebuit

en smeekte om gratie.
De tegenprestatie
bleek een penetratie.

Zo was de vertelling
in haar boerengezin
na Emmy’s bevrijding.

Het duel was beslecht.
Emmy won het gevecht.
Heidi was afgelegd.

Ze staakte haar verweer
en boog haar hoofd diep neer.
De jager wilde meer.

Hij trok aan de ketting
van de strijdkoningin.
Toen ze naar hem toe ging,

heulend met zijn acties,
bleek zijn wens onethisch.
Strak tegen de tralies

plette hij haar lichaam.
Hoewel onaangenaam
verdroeg ze het lijdzaam.

Al haar ledematen
die ineens vastzaten
via stalen platen

aan de deur van de cel,
waren in de knel.
Haar heer speelde vals spel.

De verstoorde balans
gaf Heidi toen de kans.
De weerwraak stimulans

kon Heidi niet weerstaan.
Ze liet ze totaal gaan.
Wat haar was aangegaan

moest Emmy ervaren.
Ze trok aan schaamharen
tot ze eraf waren.

Ze trok ze allemaal.
Emmy’s kutje werd kaal.
Ze maakte geen kabaal

hoewel haar mond scheef trok
als weer een pluk vertrok.
Emmy was gans de bok.

De jager spoorde aan
om vooral door te gaan
tot alles was gedaan.

Dat had hij veel liever.
Ze leek dan naïever
en bleek effectiever

om de besmeurde pruim
te ontdoen van melkschuim
na diepgaand plichtsverzuim.

Heidi’s vergeldingsdrang
bleek nog niet in bedwang.
Emmy haar rechterwang

kreeg met de hand als roe.
Ze draaide uit armoe
ook haar linkerwang toe

om de cyclus van haat
met een tedere daad
te doorbreken als kwaad.

Met deze symboliek
als zwijgende repliek
verstomde de tragiek.

De vuige jager zei;
“de winnaar is Heidi.
Zij krijgt mijn provisie.”

Zo begon het vozen.
“Hij heeft haar gekozen
om zijn zaad in te lozen

en verkoopt het als prijs.
Mij maak je meer niet wijs
dat je jouw levensreis

slechts met mij wil delen.
Je koos voor valsspelen
om haar te gaan strelen

en in haar te douwen.
Je bent dus met vrouwen
geenszins te vertrouwen.

Als meisje met fatsoen
zal ik je gaan voordoen
het juiste te gaan doen.

Ik zal iets verzinnen
je voor me te winnen,
je laten bezinnen

en laten berouwen.
Je stopt dan met vrouwen
en gaat met mij trouwen.”,

zo mijmerde Emmy
bij de expositie
van Emma’s conceptie.

Heidi’s opdrachtgever
dronk pruimenjenever
tijdens zijn gezever

over bedprestaties.
De dronkaard vond Heidi’s
nogal koud en klinisch,

verzuimend dat een crème
hem hielp bij zijn handrem.
Daarnaast lag Emma klem

geboeid en afgekneld
door seksueel geweld.
Toen Emma was geveld

droop de spons in en af
terwijl hij overgaf.
Hij bedacht nog een straf

voor hij beiden verliet.
De bespuugde pisgriet
wilde eerst liever niet,

maar moest naar Emmy gaan
en tegen haar aan staan.
Ze raakte elkaar aan;

de kutjes, de tieten,
van de naakte grieten.
Hij kon wel genieten

van de dichte luikjes
onder ieders pruikjes
toen hij beide buikjes

met zijn broekriem verbond.
Van schaamte door de grond
was de pijn die hij verslond.

Beschroomdheid door weerzin.
Schuchterheid door walging.
Schroom door vernedering.

“Hou dat twistvuur brandend
al voelt het aanrandend!”,
zei hij klappertandend;

“Doe het uit lijfsbehoud
anders wordt het te koud.”
Hij gooide wat brandhout

in een koperen stoof
die hij naar hen toe schoof.
Emmy had ongeloof

of dit genoeg bijdroeg.
De avond was nog vroeg.
Het hout was niet genoeg.

Hij liet het tafereel,
verdween van het toneel
en liep naar zijn kasteel.
